% LARMIP contribution: ACCESS-ISSM (AIS3 configuration) -- DRAFT 2026-10-01.
% Before sending: confirm the model/version name (CTRL drift filled in from the 200-yr run), and
% fix the three TODO entries in larmip_refs.bib.
% Uses natbib (\citet/\citep) with larmip_refs.bib.

\paragraph{ACCESS-ISSM (AIS3).}
The simulations use the Ice-sheet and Sea-level System Model \citep[ISSM;][]{larour2012} through
its Python interface, pyISSM, in a continental Antarctic configuration developed for ACCESS. Ice
flow is solved with the shallow-shelf approximation \citep{macayeal1989} on an unstructured
triangular mesh of 1.59 million vertices, refined by observed speed: element edges are at most
1.5\,km on fast-flowing grounded ice, 0.5--5\,km on ice shelves, and the median edge length over the
ice sheet is 1.5\,km. Bed topography and initial ice geometry are from BedMachine Antarctica v3
\citep{morlighem2020}. Basal sliding follows a Budd-type law \citep{budd1979} with linear dependence
on velocity and effective pressure ($p = q = 1$). Effective pressure is taken from the subglacial
hydrology estimate of \citet{ehrenfeucht2024} and bounded below by 7\,\% of the ice overburden
pressure. The flow law uses $n = 3$. On grounded ice the rate factor follows from RACMO2.4p1
surface temperature \citep{vandalum2024} via the temperature dependence of \citet{cuffey2010}. On
floating ice it is obtained by inversion. The grounding line evolves with sub-element
parameterisations of its position and of basal friction \citep{seroussi2014}; no basal melt is
applied to partially floating elements. The calving-front position is held fixed, so the
experiments isolate the response to melt. Sub-shelf melt
uses the ISMIP6 local quadratic parameterisation \citep{jourdain2020} driven by an observational
thermal-forcing climatology \citep{zhou_tf}. Its coefficient ($\gamma_0 = 300$\,m\,yr$^{-1}$) and
per-basin temperature corrections were calibrated, following the ISMIP7 ocean protocol, to
basin-integrated present-day melt \citep{paolo2023,adusumilli2020} in IMBIE2 basins
\citep{mouginot2017}, with buttressing-weighted skill \citep{reese2018} used to choose $\gamma_0$.

\paragraph{Initialisation and spin-up.}
The model is initialised by data assimilation rather than a long thermal spin-up. The basal
friction coefficient is inferred by minimising the misfit to MEaSUREs InSAR-based surface
velocities \citep{rignot2011,mouginot2019} with an adjoint method \citep{morlighem2010} and the
M1QN3 optimiser \citep{gilbert1989}. The cost function combines absolute and logarithmic velocity
misfit with a Tikhonov regularisation of the friction gradient; the resulting grounded-ice velocity
RMSE is 66\,m\,yr$^{-1}$ in a full-continent forward solve. From this state the model is run from
1995 to 2019 with annual RACMO2.4p1 surface mass balance and the calibrated sub-shelf melt;
modelled elevation change over that period was compared with CryoSat-2-era altimetry. The 2019 end
state is the starting point for all LARMIP experiments. All experiments hold surface mass balance
at its 1995--2014 RACMO2.4p1 mean and the ocean forcing constant. Over the 200 years the control
gains ice at a slowly decreasing rate, from 0.56 to 0.39\,mm\,yr$^{-1}$ of sea-level equivalent
($-89$\,mm in total), with mass loss only in the Amundsen Sea sector. All results are reported relative to the control.
In the LARMIP experiments the prescribed additional basal melt is added to the parameterised melt
under every fully floating element in the respective region, so the forced area follows the
grounding line as it migrates.
